\scanpage{45}

a comme anneau de coeff. $W(k)$, dont le corps des fractions $K$ est de car. 0, et non $k$. De plus, cette approche \deleted{ne} \inserted{(s'il est vrai qu'elle marche sur $W(k)$, pas seulement $K$)} \deleted{est} fournit \deleted{des relations satisfaisantes} une notion de torsion $p$-adique, \deleted{entre la cohomologie} qui est peut-être la bonne : est-ce la torsion $p$-adique de la fibre générique\deleted{?} géométrique ? Ici est \unsure{resoulevée} la question, extrêmement importante et profonde, des relations entre la cohomologie de D.R. et la cohomologie $p$-adique, quand on spécialise de la car. 0 à la car. $p>0$. (question déjà envisagée par Tate, \inserted{dans un langage différent,} dans son \unsure{séminaire} l'an dernier dans le \inserted{cas des} $H^1$ tout au moins).

Sur l'importance de cette question \deleted{\unsure{et}} pour l'attaque de certains pbs liés aux \inserted{familles de} cycles algébriques, cf mes papiers aux Inventiones \inserted{$[\,]$} à paraître prochainement.

\margin{44}
